\begin{textsample}{NEG Dim 2 – 2019\_c5 – Score -3.00 – 2019\_c5/enade\_2019\_enf\_ba60297fd1.txt}  \label{ex:f2_neg_093}
CAPÍTULO 5 DISTRIBUIÇÃO DOS CONCEITOS 5.1 PANORAMA NACIONAL DA DISTRIBUIÇÃO DOS CONCEITOS O Conceito Enade é um indicador de qualidade que avalia os cursos por intermédio dos desempenhos dos estudantes a partir dos resultados do Enade.
Os valores possíveis vão de 1 ( pior situação ) a 5 ( melhor situação ).
Os cursos com apenas um ou sem nenhum concluinte participante, não obtêm o Conceito Enade, ficando Sem Conceito ( SC ).
Na Tabela 5.1, são apresentadas a quantidade e a distribuição de cursos de Enfermagem participantes do Enade/2019, por faixa de conceito e Grande Região. \textbf{Enfatiza} -se, mais uma vez, que a diferença entre os cursos tabulados neste capítulo e no capítulo 2 corresponde aos cursos que não foram avaliados, ou seja, àqueles, em princípio, sem alunos concluintes que tivessem sido inscritos para a prova.
Observando -se os dados da Tabela 5.1, nota -se que, dos 800 cursos participantes, 310 ( 38,8 \% ) classificaram -se com conceito 3, o valor modal.
Esse foi, também, o conceito modal nas regiões : Sudeste ( 42,9 \% ) e Sul ( 56,4 \% ).
O conceito 2 foi o segundo mais frequente em nível nacional ( 33,6 \%, correspondendo a 269 cursos ), sendo o conceito modal da maioria das regiões : Norte ( 50,0 \% ), Nordeste ( 38,3 \% ) e Centro-Oeste ( 37,2 \% ).
Por sua vez, o conceito 4 foi o terceiro mais frequente ( 16,4 \%, correspondendo a 131 cursos ).
Houve, ainda, 54 cursos ( 6,8 \% ) que receberam conceito 1, e 29 cursos ( 3,6 \% ) que receberam conceito 5.
Dos 800 cursos de Enfermagem, sete cursos ( 0,9 \% ) ficaram Sem Conceito ( SC ).
A região Norte participou com 62 cursos ou 7,8 \% do total nacional.
Desses, 31 cursos receberam o conceito 2, o que equivale a 50,0 \% do total regional, além de ser o conceito modal na região, como já comentado.
A nove cursos ( 14,5 \% ) atribuiu -se conceito 1, e a 14 cursos ( 22,6 \% ), conceito 3.
Os conceitos 4 e 5 foram atribuídos a seis e um curso, respectivamente, sendo 9,7 \% e 1,6 \%.
Um curso ficou Sem Conceito ( SC ) na região Norte.
A região Nordeste participou com 209 cursos ou 26,1 \% do total nacional.
Nessa região, 80 cursos ( 38,3 \% em termos regionais ) obtiveram conceito 2, o conceito modal para a região, como já comentado.
Os conceitos 1 e 3 foram atribuídos, respectivamente, a 17 e a 60 cursos ( 8,1 \% e 28,7 \%, respectivamente ).
Ainda nessa região, 36 cursos ( 17,2 \% ) foram avaliados com conceito 4, e 14 cursos ( 6,7 \% ), com o conceito 5.
Dois cursos ficaram Sem Conceito ( SC ).
Dos 326 cursos participantes da região Sudeste ( 40,8 \% do total nacional ), o conceito 3 foi o conceito modal, como já anteriormente assinalado, tendo sido obtido por 140 cursos ( 42,9 \% ).
O conceito 1 foi atribuído a 17 cursos ( 5,2 \% ), e o conceito 2, a 104 cursos ( 31,9 \% ).
Receberam os conceitos 4 e 5, respectivamente, 55 cursos ( 16,9 \% ) e sete cursos ( 2,1 \% ).
Nessa região, três cursos ficaram Sem Conceito ( SC ).
Os 117 cursos da região Sul corresponderam a 14,6 \% do total nacional.
Como já indicado, houve predominância do conceito 3, atribuído a 56,4 \% dos cursos da região ( conceito modal ), ou seja, atribuído a 66 dos 117 cursos participantes na região Sul.
O conceito 1 foi atribuído a três cursos ( 2,6 \% ), os conceitos 2 e 4, a 22 cursos, cada ( 18,8 \% cada ), e o conceito 5 foi atribuído a quatro cursos ( 3,4 \% ).
Na região Sul, nenhum curso ficou Sem Conceito ( SC ).
Já dos 86 cursos participantes na região Centro-Oeste ( 10,8 \% do total nacional ), 32 ( 37,2 \% em termos regionais ) receberam conceito 2, o conceito modal, como já fora destacado.
Com uma concentração próxima a modal, o conceito 3 foi atribuído a 30 cursos ( 34,9 \% ).
Oito cursos ( 9,3 \% ) obtiveram conceito 1, 12 ( 14,0 \% ), o conceito 4 e três ( 3,5 \% ), conceito 5.
Além disso, um curso ficou Sem Conceito ( SC ).
No Gráfico 5.1, é apresentada a distribuição cumulativa do Conceito Enade segundo Grande Região.
Em um gráfico de distribuição cumulativa, em cada valor de abscissa ( eixo x ), o valor apresentado no eixo das ordenadas ( eixo y ) é o valor acumulado da variável ( conceito Enade ).
Observando a linha vertical pontilhada, que ajuda a visualizar o percentual de conceitos acumulados até o valor 3, ou seja, conceitos 1, 2 e 3, pode -se notar que, por exemplo, a região Norte ( linha azul ) apresenta quase 90 \% dos cursos com conceito 3 ou abaixo deste valor.
A região Nordeste ( linha verde ), por outro lado, apresenta cerca de 75 \% dos cursos na mesma situação : com conceito 3 ou abaixo.
Como em todas as regiões houve cursos com conceito 5, as poligonais cumulativas só alcançam 100 \% no conceito máximo.
Em linhas gerais, em um gráfico de distribuição cumulativa, poligonais mais à esquerda representam Grandes Regiões com pior distribuição, e poligonais mais à direita, Grandes Regiões com melhores conceitos.
A região Norte ( linha azul ) apresenta os cursos com conceitos mais baixos, ao passo que a região Sul ( linha rosa ) apresenta as distribuições com valores maiores.
Cursos em IES situadas nas regiões Nordeste ( linha verde ), Sudeste ( linha vermelha ) e Centro-Oeste ( linha laranja ) apresentam uma situação intermediária. 5.2 CONCEITOS POR CATEGORIA ADMINISTRATIVA, POR MODALIDADE DE ENSINO E GRANDE REGIÃO Na Tabela 5.2, é apresentada a distribuição dos cursos de Enfermagem participantes do Enade/2019, por Categoria Administrativa e Modalidade de ensino, de acordo com os conceitos por eles alcançados, segundo as Grandes Regiões brasileiras.
Dos 800 cursos participantes, 142 ( 17,8 \% ) eram ministrados em Instituições Públicas, e 658 ( 82,3 \% ), em Instituições Privadas.
Quanto à Modalidade de ensino, quatro cursos eram a Distância, e 796, Presenciais.
De acordo com as informações da Tabela 5.2, em termos nacionais, entre Instituições Públicas e Privadas, observa -se que, dos 29 cursos avaliados com conceito 5, 26 eram oferecidos em IES Públicas, e três, em IES Privadas.
Dos 142 cursos participantes de IES Públicas, o conceito 4 foi o valor modal, atribuído a 71 cursos ( 50,0 \% ).
Dos demais cursos de Instituições Públicas participantes, três cursos ( 2,1 \% ) foram avaliados com conceito 1, 14 cursos ( 9,9 \% ) foram avaliados com conceito 2, 27 cursos ( 19,0 \% ) foram avaliados com conceito 3 e, como já comentado, 26 cursos ( 18,3 \% ), com conceito 5.
Um curso ficou Sem Conceito ( SC ).
Na rede privada, o conceito modal foi 3, atribuído a 283 cursos, o correspondente a 43,0 \% dos 658 cursos da categoria.
Dos demais cursos participantes, 51 ( 7,8 \% ) receberam conceito 1 e 255 ( 38,8 \% ), conceito 2.
O conceito 4 foi atribuído a 60 cursos ( 9,1 \% ), e o conceito 5, a três cursos ( 0,5 \% ).
Nesta Categoria Administrativa, seis cursos ( 0,9 \% ) ficaram Sem Conceito ( SC ).
Dos 796 cursos participantes da modalidade presencial, o conceito 3 foi o valor modal, atribuído a 310 cursos ( 38,9 \% ).
Dos demais cursos nesta modalidade de ensino, 52 cursos ( 6,5 \% ) receberam conceito 1, 267 cursos ( 33,5 \% ) receberam conceito 2, 131 cursos ( 16,5 \% ), conceito 4, e 29 cursos ( 3,6 \% ), conceito 5.
Nesta modalidade de ensino, sete cursos ( 0,9 \% ) ficaram Sem Conceito ( SC ).
Aos quatro cursos a Distância informados atribuíram -se os conceitos 1 e 2, dois cursos cada.
Na análise por região, observa -se que, na região Norte, as Instituições Privadas participaram com 48 cursos ( 77,4 \% do total regional ), dos quais nenhum ficou Sem Conceito ( SC ) ou recebeu o conceito 5.
O conceito modal para as Instituições Privadas na região foi 2, atribuído a 29 cursos, correspondendo a 60,4 \% dessa categoria na região.
Nove cursos foram avaliados com conceito 1, outros nove cursos foram avaliados com conceito 3, e um curso, com conceito 4.
As Instituições Públicas participaram com 14 cursos na região Norte ( 22,6 \% do total regional ), aos quais foram atribuídos os conceitos 2 ( dois cursos ), 3 ( cinco cursos ), 4 ( outros cinco cursos ) e 5 ( um curso ).
Um curso ficou Sem Conceito ( SC ) e nenhum com conceito 1 nessa combinação de categoria e região.
Todos os cursos oferecidos na região Norte eram presenciais.
Na região Nordeste, a rede privada concentrou 153 dos 209 cursos participantes, o equivalente a 73,7 \% do total da região.
Dentre os cursos oferecidos por IES Privadas no Nordeste, 77 cursos foram avaliados com conceito 2, o conceito modal.
Dois cursos ficaram Sem Conceito ( SC ), 16 receberam o conceito 1, 52, o conceito 3, seis, o conceito 4, e um curso recebeu o conceito 5.
As Instituições Públicas da região Nordeste participaram com 55 cursos ( 26,3 \% do total da região ).
Desses, 13 obtiveram conceito 5, 30, o conceito 4, oito o conceito 3, três, o conceito 2, e um curso obteve o conceito 1.
Nenhum curso ficou Sem Conceito ( SC ).
Todos os cursos oferecidos na região Nordeste eram presenciais.
Na região Sudeste, a proporção de cursos da rede privada, 90,8 \%, foi mais \textbf{elevada} do que a registrada nas demais regiões brasileiras, correspondendo a 296 dos 326 cursos participantes.
Nessa categoria e região, o conceito modal foi 3, atribuído a 136 cursos, e três cursos ficaram Sem Conceito ( SC ).
Os demais foram avaliados com conceito 1 ( 17 cursos ), conceito 2 ( 100 cursos ), conceito 4 ( 39 cursos ) e conceito 5 ( um curso ).
Dos 30 cursos oferecidos por Instituições Públicas na região Sudeste, o conceito modal foi 4, atribuído a 16 cursos.
Nessa região, nenhum curso de IES Públicas ficou Sem Conceito ( SC ) ou recebeu conceito 1.
Quatro cursos receberam conceito 2, outros quatro o conceito 3, e seis cursos receberam o conceito 5.
Também na região Sudeste, a maioria dos cursos foi oferecida na modalidade Presencial : 324 cursos.
Para esses, a moda foi o conceito 3, com 140 cursos.
Os demais cursos presenciais da região foram alocados aos conceitos 1 ( 17 cursos ), 2 ( 102 cursos ), 4 ( 55 cursos ) e 5 ( sete cursos ).
Ficaram Sem Conceito ( SC ) três cursos.
Na região Sudeste, os dois cursos oferecidos na modalidade Educação a Distância obtiveram o conceito 2.
As Instituições Privadas concentraram 95 dos 117 cursos participantes da região Sul, 81,2 \% do total regional.
Desses, 59 cursos obtiveram conceito 3, o conceito modal.
Nesta combinação de Categoria Administrativa e Grande Região, dois cursos receberam conceito 1, 22 cursos receberam conceito 2, 11 cursos, o conceito 4, e ainda um, conceito 5.
Nenhum curso ficou Sem Conceito ( SC ).
As Instituições Públicas na região Sul participaram com 22 cursos ( 18,8 \% dos cursos da região ), aos quais foram atribuídos os conceitos 1 ( um curso ), 3 ( sete cursos ), 4 ( 11 cursos, modal ) e 5 ( três cursos ).
Nenhum curso da região Sul oferecido por IES Pública ficou Sem Conceito ( SC ) ou recebeu o conceito 2.
Todos os cursos oferecidos na região Sul eram presenciais.
Na região Centro-Oeste, 65 dos 86 cursos participantes eram de Instituições Privadas ( 75,6 \% em termos regionais ).
A distribuição dos cursos pelos conceitos nessa combinação de categoria e região foi bimodal, com 27 cursos nos conceitos 2 e 3, cada.
O conceito 1 foi atribuído a sete cursos, e o conceito 4, a três cursos.
Um curso ficou Sem Conceito ( SC ) e nenhum curso recebeu conceito 5.
Quanto aos 21 cursos oferecidos por Instituições Públicas na região Centro-Oeste ( 24,4 \% do total regional ), nove cursos apresentaram conceito 4, conceito modal.
Os demais cursos foram \textbf{distribuídos} por todos os conceitos, sendo : 1 ( um curso ), 2 ( cinco cursos ), 3 ( três cursos ) e 5 ( mais três cursos ).
Nenhum curso ficou Sem Conceito ( SC ) nessa combinação de categoria e região.
A maioria dos cursos oferecidos na região Centro-Oeste eram presenciais, apenas dois cursos eram a Distância e obtiveram conceito 1.
Dos 84 cursos presenciais dessa região, a moda foi o conceito 2, com 32 cursos, seguido de perto pelo conceito 3, com 30 cursos.
Os demais cursos foram alocados aos seguintes conceitos : seis cursos ( conceito 1 ), 12 cursos ( conceito 4 ) e três cursos ( conceito 5 ).
Nessa combinação de modalidade e região, um curso ficou Sem Conceito ( SC ).
No Gráfico 5.2, apresenta -se a distribuição do Conceito Enade, segundo a Categoria Administrativa da IES.
Os cursos em IES Públicas ( linha azul ) apresentam uma poligonal mais à direita do que os em IES Privadas ( linha verde ) e, consequentemente, uma distribuição de Conceitos Enade com valores maiores.
Como o conceito 5 não foi alcançado por nenhum curso de IES Privada a poligonal correspondente ( linha verde ) alcança 100 \% no conceito 4. 5.3 CONCEITOS POR ORGANIZAÇÃO ACADÊMICA E GRANDE REGIÃO Na Tabela 5.3, encontra -se a distribuição dos conceitos atribuídos aos cursos participantes do Enade/2019, na Área de Enfermagem, por Organização Acadêmica, segundo as Grandes Regiões brasileiras.
Dos 800 cursos de Enfermagem participantes, 292 eram oferecidos em Universidades, 196 em Centros Universitários, 310, em Faculdades e dois, em CEFET/IF.
Esta distribuição corresponde a, respectivamente, 36,5 \%, 24,5 \%, 38,8 \% e 0,3 \% dos cursos.
De acordo com os dados apresentados, de todos os 29 cursos avaliados com conceito 5, 24 eram vinculados a Universidades.
Os cursos vinculados a esse tipo de Organização Acadêmica tiveram o conceito 3 como conceito modal, atribuído a 108 cursos ( 37,0 \% ).
Os demais cursos vinculados a Universidades receberam os conceitos 1 ( cinco cursos ), 2 ( 55 cursos ), 4 ( 98 cursos ) e conceito 5 ( 24 cursos, como já mencionado ).
Dois cursos ficaram Sem Conceito ( SC ).
Dos cursos em Centros Universitários, o conceito modal também foi 3, atribuído a 92 cursos ( 46,9 \% ).
Nenhum curso vinculado a esse tipo de Organização Acadêmica ficou Sem Conceito ( SC ), e os outros cursos receberam os conceitos 1 ( oito cursos ), 2 ( 74 cursos ), 4 ( 20 cursos ) e 5 ( dois cursos ).
Dos 310 cursos mantidos por Faculdades, 140 ( 45,2 \% ) receberam conceito 2, o conceito modal.
O segundo conceito mais frequentemente atribuído aos cursos oferecidos por essa categoria de IES foi o conceito 3, recebido por 109 cursos.
Na sequência, 41 cursos obtiveram conceito 1, 12 cursos, conceito 4, e três cursos, conceito 5.
Dos cursos oferecidos em Faculdades, cinco ficaram Sem Conceito ( SC ).
Os dois cursos mantidos por CEFET/IF, foram \textbf{distribuídos} igualmente aos conceitos 3 e 4.
Considerando -se, separadamente, as regiões brasileiras, verifica -se que, na região Norte, as Universidades foram representados por 17 cursos.
O conceito 5 foi atribuído a um curso, os conceitos 2, 3 e 4 foram atribuídos a cinco cursos cada.
Um curso ficou Sem Conceito ( SC ) e nenhum recebeu o conceito 1 nesta combinação de categoria e região.
Os Centros Universitários da região Norte participaram com, apenas, 12 cursos avaliados, onde o conceito modal foi 3 concentrando seis cursos e seguido de perto pelo conceito 2 com cinco cursos.
Um curso recebeu conceito 4, nenhum curso ficou Sem Conceito ( SC ) ou recebeu os conceitos 1 ou 5.
As Faculdades participaram com 33 cursos na região Norte, onde o conceito modal foi 2 com 21 cursos.
O conceito 1 foi atribuído a nove cursos, e o conceito 3, a três cursos, além de nenhum curso ter ficado Sem Conceito ( SC ) e, também, não houve cursos avaliados com os conceitos 4 ou 5.
Não foram avaliados cursos na categoria CEFET/IF na região Norte.
Na região Nordeste, as Universidades participaram com 67 dos 209 cursos da Área de Enfermagem oferecidos.
Nenhum curso ficou Sem Conceito ( SC ) ou foi avaliado com o conceito 1.
O conceito modal 4 foi atribuído a 31 cursos.
Os cursos restantes obtiveram conceitos 2 ( 11 cursos ), 3 ( 12 cursos ) e 5 ( 13 cursos ).
Os Centros Universitários contaram com 49 cursos participantes na região Nordeste, 24 dos quais receberam o conceito modal 2.
Para os demais cursos, foram atribuídos os conceitos 1 ( um curso ), 3 ( 21 cursos ), 4 ( dois cursos ) e 5 ( um curso ).
Nenhum curso oferecido por Centros Universitários no Nordeste ficou Sem Conceito ( SC ).
As Faculdades foram representadas por 92 cursos na região Nordeste, 45 dos quais receberam conceito 2, o conceito modal.
A mesma quantidade de cursos do conceito modal foi alocada aos conceitos 1 ( 16 cursos ), 3 ( 27 cursos ) e 4 ( dois cursos ).
Na região Nordeste, dois cursos mantidos por Faculdades ficaram Sem Conceito ( SC ), e nenhum curso recebeu o conceito 5.
Os CEFET/IF participaram com um curso na região Nordeste avaliado com conceito 4.
Na região Sudeste, as Universidades concentraram 119 dos 326 cursos de Enfermagem da região.
Dos cursos oferecidos em Universidades, nessa região, o conceito modal foi 3, atribuído a 52 cursos.
Os demais cursos receberam os conceitos 1 ( três cursos ), 2 ( 25 cursos ), 4 ( 34 cursos ) e 5 ( quatro cursos ).
Um curso ficou Sem Conceito ( SC ).
Os Centros Universitários participaram com 88 cursos na região Sudeste, dos quais 41 obtiveram o conceito modal 3.
O conceito 1 foi atribuído a seis cursos, o conceito 2, a 29 cursos, o conceito 4, a 11 cursos e o conceito 5, a um curso.
Nenhum curso ficou Sem Conceito ( SC ) nessa combinação de organização acadêmica e região.
Já as Faculdades foram representadas por 119 cursos na região Sudeste, dos quais dois ficaram Sem Conceito ( SC ).
O conceito modal 2 foi atribuído a 50 cursos.
Aos demais cursos foram atribuídos os conceitos 1 ( oito cursos ), 3 ( 47 cursos ), 4 ( dez cursos ) e 5 ( dois cursos ).
Não foram avaliados cursos na categoria CEFET/IF na região Sudeste.
Dos 117 cursos da região Sul, 61 eram mantidos por Universidades, para os quais o conceito modal foi 3, atribuído a 30 cursos.
Os demais cursos receberam os conceitos 2 ( oito cursos ), 4 ( 19 cursos ) e 5 ( quatro cursos ).
Nesse tipo de organização, nenhum curso da região Sul ficou Sem Conceito ( SC ) ou recebeu conceito 1.
Os 25 cursos participantes de Centros Universitários da região Sul receberam os conceitos intermediários 2 ( seis cursos ), 3 ( 16 cursos, modal ) e 4 ( três cursos ).
Nenhum curso ficou Sem Conceito ( SC ) ou recebeu os conceitos extremos 1 ou 5.
Quanto aos 30 cursos vinculados a Faculdades na região Sul, 19 receberam o conceito modal 3.
O conceito 2 foi atribuído a oito cursos, e o conceito 1, a três cursos.
Nenhum curso ficou Sem Conceito ( SC ) ou recebeu os conceitos mais altos 4 ou 5.
Já os CEFET/IF participaram com um curso na região, ao qual foi atribuído conceito 3.
Na região Centro-Oeste, 28 dos 86 cursos oferecidos eram mantidos por Universidades.
Para este tipo de Organização Acadêmica, nenhum curso ficou Sem Conceito ( SC ).
Os cursos receberam os conceitos : 1 ( dois cursos ), 2 ( seis cursos ), 3 e 4 ( nove cursos cada, bimodal ) e conceito 5 ( dois cursos ).
Os Centros Universitários da região Centro-Oeste contaram com 22 cursos, todos avaliados com conceito, sendo 1 curso, com conceito 1, dez cursos, com conceito 2 ( modal ), oito cursos, com conceito 3, e quatro cursos, com conceito 4.
Dos 36 cursos oferecidos em Faculdades na região Centro-Oeste, 16 receberam o conceito modal 2, cinco receberam o conceito 1, 13, o conceito 3, e um, o conceito 5.
Nesta região, um curso oferecido em Faculdade ficou Sem Conceito ( SC ), e nenhum obteve conceito 4.
Não foram avaliados cursos na categoria CEFET/IF na região Centro-Oeste.
No Gráfico 5.3, apresenta -se a distribuição do Conceito Enade dos cursos de Enfermagem, segundo a Organização Acadêmica da IES.
Os cursos em CEFET/IF ( linha laranja ) apresentam uma poligonal mais à direita que as demais, denotando uma melhor distribuição de Conceitos Enade.
No outro extremo, as Faculdades ( linha vermelha ), apresentando a respectiva poligonal mais à esquerda e, consequentemente, distribuição de Conceitos Enade com menores valores.
As Universidades ( linha azul ) e Centros Universitários ( linha verde ) apresentam um comportamento intermediário.
Como no Gráfico 5.2 pode -se observar que a poligonal relativa aos CEFET/IF ( linha laranja ) alcança 100 \% no conceito 4 e isso ocorre porque nenhum curso dessa Organização Acadêmica recebeu conceito 5.

% matched lemmas: distribuir, elevado, enfatizar
\end{textsample}
